% Méthode proposée : la réalisation P est déjà définie avant cet encodage.
\begin{frame}[t]{Encoder un polyèdre : pseudo-code}
\small
\textbf{Entrées :} primitives $\mathcal F\ne\varnothing$ ; grille ordonnée $B=(b_j)_{j=1}^{512}$.
\par
\vspace{.23cm}
\begin{beamercolorbox}[sep=9pt,wd=\textwidth]{block body}
\renewcommand{\arraystretch}{1.35}
\begin{tabular}{@{}r@{\hspace{.28cm}}l@{}}
1 & $c\gets$ centre de la boîte englobante de $\mathcal F$\\
2 & $s\gets$ moitié de son plus grand côté\\
3 & \textbf{si} $s=0$ : \textbf{retourner} $\bigl((\norm{b_j})_j,\ c,\ 0\bigr)$\quad\textcolor{gray}{(un point)}\\
4 & $\widehat{\mathcal F}\gets\{(\sigma-c)/s:\sigma\in\mathcal F\}$\\
5 & \textbf{pour chaque} sonde $b_j$, dans l'ordre fixé :\\
6 & \hspace{.6cm}$D[j]\gets\displaystyle\min_{\widehat\sigma\in\widehat{\mathcal F}}\mathrm{distance}(b_j,\widehat\sigma)$\\
7 & \textbf{retourner} $(D,c,s)$
\end{tabular}
\end{beamercolorbox}
\source{Sondes : centres des $8^3$ cellules de $[-1,1]^3$. Distance aux primitives entières, nulle à l’intérieur d’un tétraèdre plein. Attributs et relations stockés à part.}
\end{frame}
